\begin{abstract}
Chain-of-thought (CoT) prompting improves large language model (LLM) reasoning by eliciting intermediate steps before a final answer. Standard CoT uses formulaic step-by-step decomposition, but cognitive science suggests humans frequently reason through narrative mental simulation---constructing stories with characters, settings, and causal progressions. We introduce \storycot, a prompting strategy in which few-shot exemplars embed the same logical steps within vivid narratives, and conduct the first systematic comparison of narrative-form CoT against standard CoT across four diverse reasoning benchmarks: \gsm (math), \csqa (commonsense), \strategyqa (multi-hop), and \arc (science). Using \gptfour with 150 questions per dataset, we find that \storycot performs \emph{equivalently} to standard \fewshotcot across all domains, with accuracy differences of 0--1.3 percentage points and no statistically significant differences (McNemar's test, all $p > 0.47$). Both methods dramatically outperform direct prompting on tasks requiring multi-step reasoning (up to +36 pp on \gsm). Our results suggest that for state-of-the-art LLMs, the benefit of CoT arises from computational scaffolding---providing intermediate token positions for reasoning---rather than from any particular reasoning format. This finding has practical implications: narrative-form prompts can be used without performance penalty in applications where natural-sounding reasoning is preferred.
\end{abstract}
